% Transcription — fonds Alexandre Grothendieck, shelfmark 15, batch 6 (pages 101-120).
% Established from the facsimile archives/batches/15/batch-06.pdf (a mirror of
% https://grothendieck.umontpellier.fr/15.pdf), per the skill
% transcribe-grothendieck. The batch file carries no cover sheet: PDF page N
% is page 100+N of the folder (the Montpellier footer prints « 101 » ...
% « 120 »), and the archivists' pencil numbers bottom-left agree on every
% transcribed leaf (101, 103, 105, 107, 109, 111, 113, 116, 118, 120 read).
%
% Pass: Opus 5.5 (claude-opus-5-5), 2026-09-23 — first pass, unchecked against the pages by a human.
%
% Ten of the twenty pages are transcribed: the odd pages 101 to 113, then
% 116, 118, 120. Absent: page 115, blank; the others are versos of unrelated
% typescripts, nothing of his mathematics on them — 102, an English
% typescript « IV.8 » on ringed spaces and modules (φ*, φ_*, canonical
% homomorphisms); 104, 106, 110, 112, 114, 119, the English translation typed
% in violet already met in batch 5 (typed pp. 94, 92, 86, 84, 85, 78:
% sections of line bundles, biregular embeddings, induction on dimension;
% margin « p. 213 » typed on 112); 108, a typescript on maximal tori and
% unipotent elements (« 2.9 Corollary », « 2.10 Remark »), upside down; 117,
% typed errata to an exposé on groups of multiplicative type (« Page 33.
% Ligne -9 à -1, remplacer le texte par le suivant : Remarques 7.10 ... »),
% with one line in his hand (« Page 35, ligne 9, 10, supprimer "lisse sur
% k" ») which concerns that typescript only. All skipped per the project's
% rule on typed versos.
%
% What the batch is. One run, in his hand, blue-black ink, the rectos
% continuing batch 5's pages 83-99 (not paginated by him): the motivic Galois
% group of finite fields, Weil q-numbers, CM fields, and the real fibre
% functor.
%   101    continues « Vérifions que cela donne bien » (p. 99): G/G° is dual
%          to a group of roots of unity, whose dual is that of Ẑ; M(q) ⊂ Q'(i)*
%          the Weil q-numbers, G(q) its dual, f_q ∈ G(q), the motivic Galois
%          group of F_q; M(q^m) ⊂ M(q), M(q)^m ⊂ M(q^m); G(q) ← G(q^m),
%          (f_q)^m ↤ f_{q^m}; the transpose is raising to the m-th power;
%          margin: M(q) ≃ Z × M°(q), M° the units exc. p with
%          v_p(u) + v_p'(u) = 0; inductive limit M(q^∞).
%   103    M(q^∞) ≃ Z × M°(q^∞), M°(q^∞) = lim M°(q^m) ≃ M° ⊗_Z Q, which is a
%          quotient P/U (U roots of unity, P the λ ∈ Q̄* with λ^n ∈ M°(q));
%          a boxed attempt at M°(q = p^n) → P/U, λ ↦ λ^{1/n} mod U, cancelled
%          by two long diagonals, ending on the question whether for
%          λ ∈ M°(p), n ∈ N* there is h with λ^h ∈ M°(p^{hn}).
%   105    « forme plus juste, "géométrique" » of the conjecture: the motivic
%          Galois group of F̄_p would be G ≃ D(M° ⊗ Q) × G_m, j the injection
%          of G_m, ε = (x, λ) ↦ λ²; « Conjecture affaiblie » in terms of
%          Frobenius eigenvalues λ = p^{i/2}u (λ^m eigenvalue of f_{p^m} in a
%          semi-simple motive over F_{p^m}); margin: would follow from complex
%          multiplication, cf. Taniyama-Shimura.
%   107    A) K totally real, L = K(x) quadratic: L totally imaginary ⟺
%          Δ = (x − x̄)² totally < 0; B) x ∈ Q̄: 1) x + x̄ (and x x̄) totally
%          real, (x − x̄)² totally strictly negative ⟺ 2) x generates a totally
%          imaginary quadratic L over a totally real K ⟺ 2') x ∈ Q'(i) − Q';
%          C) the same with non-strict inequality ⟺ x ∈ Q'(i) ⟺ 4) x x̄
%          totally positive for every conjugate.
%   109    diagram Q'(i), Q', L, K, L_0, K_0; M(q) the λ ∈ L with λλ̄ = q^i,
%          integral exc. p; M_0(q)/torsion ≃ (Z^P)'; a prime ℘ of L over p,
%          v_℘ : G_{m,L} → G_L; Γ_℘ the decomposition group, L_0 its fixed
%          field, K_0 = L_0 ∩ K; L_0 the field of rationality of v; L_0
%          imaginary since τ℘ ≠ ℘.
%   111    Z/2Z → Γ → Γ' → 0, Γ_℘ ⊂ Γ; (Z^P)' corresponds to a torus over Q:
%          U = Ker(∏_{L_0/K_0} G_{m,L_0} → G_{m,K_0}), G' = ∏_{K_0/Q} U, G'(Q) =
%          Ker(N : L_0* → K_0*); φ : G' → G an isogeny onto G°; v_℘ factors
%          through G'_{L_0} (diagram); f_0 lifts to G'(Q) only if M(q)/torsion
%          ≃ (Z^P)'.
%   113    the case K_0 = Q: Z/2Z acting non-trivially on Pic(A), A the
%          integers of L; it acts by symmetry, so one needs Pic(A) not a
%          2-group; margin « demander à Serre ? ». The page stops a third down.
%   116    G from a motivic situation, U = Ker(G° → G_m); Proposition 1 (U a
%          torus ⟺ ε ... ⟺ M_R has a fibre functor, 𝒢_R ≠ ∅; always true when M
%          has even degrees, j(μ_2) trivial); Proposition 2 (U = U° × μ_2 ⟺ ε
%          not a « normalised » square ⟺ 𝒢_R = ∅); proofs via an alternating
%          form and H²(R, T) for a compact torus; Cor.: the even-degree
%          subcategory M^pair admits an R-fibre functor.
%   118    Cor. in the case of Prop. 2: G° ≃ U° × G_m, ε = pr_2; the tensor
%          structure of M_R is known from G and the characteristic element of
%          H²(R, U); an object of degree 1, rank 2; when 𝒢_R ≠ ∅, H°(𝒢_R) is a
%          torsor under H¹(R, U), ₂U(R) ≃ (Z/2Z)^{dim U}; each object of 𝒢_R
%          gives a « structure polarisante » i_1 : μ_{2,C} → G°_C;
%          H°(𝒢_R) → Pol(𝒢_R) bijective.
%   120    R-fibre functors and polarisations determine each other; bracket:
%          does Aut(G°_R, ε_R, j_R) act transitively on Pol(G)? Motives over
%          F_q: a distinguished class of i : G_{m,C} → G_C with īi = j, εi = id,
%          defined by λ ↦ (1/n) v_℘(λ); « normaliser » G_C by a choice of i
%          over Q; margin example L = Q(√−p), class number 1, p = 3, 7, 11.
%          The page ends mid-sentence, « du système des conjugués », which
%          page 122 (batch 7) completes with « de i. ».
%
% Cross-batch continuations. Page 101 opens mid-sentence, completing « Cela
% détermine ... le groupe de Galois motivique des F_q. Vérifions que cela
% donne bien » at the foot of page 99 (batch 5). Page 120 ends mid-sentence
% (« par la connaissance du système des \ill{} conjugués »), continued at the
% head of page 122 (batch 7), which opens « de i. ».
%
% Notation fixed for the batch, following batch 5. \mathbf{Q}, \mathbf{Z},
% \mathbf{R}, \mathbf{C}, \mathbf{N}, \mathbf{F}_q, \mathbf{G}_m for his barred
% capitals; \mathbf{Q}' the totally real closure of Q (batch 5, p. 93),
% \mathbf{Q}'(i); \overline{\mathbf{Q}}, \overline{\mathbf{F}}_p; M(q), M°(q)
% (he also writes M_0(q), p. 109) the Weil q-numbers and their unit part;
% G(q) the dual group; \mathfrak{p} for his gothic p, \mathcal{P} for his
% script P (the set of primes of L over p, p. 109) — but P roman on page 103;
% U for his script U throughout (roots of unity on p. 103, a torus on pp. 111,
% 116, 118: different objects, as on the page); \mathcal{M} his script M (the
% motivic category, pp. 116-120); \mathcal{G}_{\mathbf{R}} his script G (the
% gerbe of real fibre functors); \Gamma, \Gamma_{\mathfrak{p}}, \Gamma'. His
% « tt », « tot. », « exc. p », « s-gpe », « homom. », « ens. », « cf. » are
% kept. Plain square brackets in the text are his.
%
% Legibility. A fast hand throughout: the formulas and statements carry, the
% connective prose of pages 101, 103, 105, 109 and 111 is often \ill{}, as are
% the slanted margins of pages 101, 116 and 120. Page 107 is written fairly
% and reads nearly end to end. Nothing on these leaves is dated.
%
% The dating is the inventory's own for the folder, copied verbatim. Its
% group, [10-18] « Théorie des motifs », is dated [à partir de 1958-à partir
% de 1982].
\documentclass[11pt,a4paper]{article}
% Le préambule commun à toutes les transcriptions du fonds Grothendieck.
%
% Il tient en une idée : **l'incertitude se compose, elle ne se résout pas.**
% Une transcription qui lisse un mot douteux a détruit la seule chose pour
% laquelle on la faisait — un lecteur doit pouvoir savoir, à chaque endroit,
% ce qui était sur le feuillet et ce qui a été ajouté. Les six macros qui
% suivent sont donc l'appareil critique, et non de la décoration.
%
% Les mêmes commandes sont reconnues par `scripts/render.mjs`, qui produit la
% vue de lecture du site. Une macro ajoutée ici doit l'être aussi là-bas,
% faute de quoi le rendu échouera bruyamment — ce qui est voulu : un rendu
% silencieusement faux est pire qu'un rendu absent.

\ProvidesPackage{grothendieck}[2026/08/08 Transcriptions du fonds Grothendieck]

\usepackage{amsmath,amssymb,amsthm}
\usepackage{mathrsfs}
\usepackage{stmaryrd}
% Les diagrammes commutatifs sont partout dans ce fonds : c'est la langue
% même de la géométrie algébrique de Grothendieck, et une transcription qui
% les rendrait en prose ne transcrirait rien.
\usepackage{tikz-cd}
% Français par défaut — c'est la langue des feuillets — mais l'anglais est
% chargé aussi : la traduction et le résumé sont des documents anglais, et la
% typographie française (espaces avant les deux-points) y serait une faute.
% Ces documents basculent par \selectlanguage{english} en tête de corps.
\usepackage[english,french]{babel}
\usepackage{fontspec}
\usepackage{geometry}
\geometry{a4paper,margin=25mm,marginparwidth=28mm}
\usepackage{marginnote}
\usepackage{xcolor}
% ulem, not soul. `soul`'s \st and \ul are built on a character-by-character
% scan that XeLaTeX's font machinery breaks: the document fails with an
% undefined control sequence rather than degrading. ulem does the same two jobs
% and is Unicode-engine safe. `normalem` stops it hijacking \emph, which the
% transcriptions use for Grothendieck's own emphasis.
\usepackage[normalem]{ulem}
\usepackage{hyperref}
\hypersetup{colorlinks=true,linkcolor=black,urlcolor=black,pdfborder={0 0 0}}

% --- Métadonnées du lot -----------------------------------------------------
% Reprises telles quelles par le script de rendu, qui les affiche en tête de
% la vue de lecture. Elles fixent ce que le fichier prétend couvrir : sans
% elles, une transcription flotte sans point d'attache dans un fonds de seize
% mille feuillets.

\newcommand{\folder}[1]{\gdef\grothFolder{#1}}
\newcommand{\batch}[1]{\gdef\grothBatch{#1}}
\newcommand{\leaves}[2]{\gdef\grothFirst{#1}\gdef\grothLast{#2}}
\newcommand{\foldertitle}[1]{\gdef\grothFoldertitle{#1}}
\gdef\grothFolder{?}\gdef\grothBatch{?}\gdef\grothFirst{?}\gdef\grothLast{?}\gdef\grothFoldertitle{}

% --- Le résumé --------------------------------------------------------------
%
% Il ouvre la lecture modernisée, sous le titre « Résumé », et il est composé
% en petit. Ce n'est pas un ornement : c'est le seul passage du document qui
% ne soit pas des mathématiques, et un lecteur doit pouvoir le reconnaître
% avant de l'avoir lu — puis le sauter s'il connaît déjà le sujet, ou s'y
% arrêter s'il ne le connaît pas. Le filet qui le suit marque la reprise du
% corps.

\newenvironment{resume}
  {\small\setlength{\parskip}{0.45em}}
  {\par\medskip\hrule\medskip}

% --- Mots-clés ---------------------------------------------------------------
%
% En anglais, en clôture du résumé de la lecture modernisée : le vocabulaire
% moderne sous lequel chercher ce que les feuillets construisent. Le manifeste
% du site (`npm run manifest`) les extrait pour étiqueter la cote entière —
% cette ligne est la source unique des tags, il n'y a pas de fichier à côté.
\newcommand{\keywords}[1]{\par\smallskip\noindent
  {\footnotesize\textsc{Keywords} --- \textit{#1}}\par}

% --- L'appareil critique ----------------------------------------------------

% Le numéro de feuillet, en marge. C'est l'ancre qui permet de régler un
% désaccord sur une formule en regardant *un* feuillet plutôt qu'une chemise
% de six cents. Le site s'en sert aussi pour tourner les pages du fac-similé
% quand on fait défiler la transcription.
\newcommand{\leaf}[1]{%
  \marginnote{\footnotesize\color{gray!70}\textsf{#1}}%
  \label{leaf:#1}\ignorespaces}

% Illisible : on ne devine pas. Un mot inventé qui se lit comme les autres est
% le pire résultat possible.
\newcommand{\ill}{\textcolor{red!60!black}{[\,\ldots\,]}}

% Lecture douteuse : le mot est donné, et signalé comme incertain.
\newcommand{\uncertain}[1]{\uline{#1}}

% Ajout de l'éditeur — une lettre manquante, un mot sous-entendu.
\newcommand{\add}[1]{\textcolor{blue!50!black}{[#1]}}

% Biffé par Grothendieck lui-même. Ce qu'il a rayé fait partie du document :
% une rature dit souvent où la pensée a changé de direction.
\newcommand{\struck}[1]{\sout{#1}}

% Note du transcripteur, sur l'état matériel du feuillet.
\newcommand{\note}[1]{\par\smallskip{\small\color{gray!60!black}\textit{[#1]}}\par\smallskip}

% Note marginale de Grothendieck, distincte du corps du texte.
\newcommand{\marginal}[1]{\par\smallskip\noindent\hspace{1em}%
  {\small\color{gray!60!black}#1}\par\smallskip}

% --- Titre ------------------------------------------------------------------

\AtBeginDocument{%
  \begin{center}
    {\large\bfseries Fonds Alexandre Grothendieck --- cote n\textsuperscript{o}~\grothFolder}\\[2pt]
    {\itshape\grothFoldertitle}\\[4pt]
    {\small feuillets \grothFirst--\grothLast\ (lot \grothBatch)}\\[2pt]
    {\footnotesize\color{gray!60!black}Transcription établie d'après le fac-similé numérisé
     par l'Université de Montpellier.\\
     Les lectures incertaines sont soulignées, les illisibles notées [\,\ldots\,],
     les ajouts entre crochets.}
  \end{center}
  \vspace{1em}}

\endinput


\folder{15}
\batch{6}
\pages{101}{120}
\dating{1965-[vers 1977]}
\watermark{Édition de démonstration}
\foldertitle{Théorie arithmétique. Théorie de Galois des motifs (1965) : notes manuscites (s.d.), tapuscrit (s.d.).}

\begin{document}

\page{101}
\note{la page s'ouvre au milieu de la phrase « Vérifions que cela donne
bien » qui termine le feuillet 99 (lot 5)}

le \ill{} correct pour les $G/G^{\circ}$ : \uncertain{celui-ci} en effet,
il \uncertain{est dual d'un} gpe des racines de l'unité, \uncertain{et ce
dernier} est bien le dual de $\hat{\mathbf{Z}}$ !

\[
  M(q^{m}) \subset M(q) \subset M(p), \qquad q = p^{n},
\]
\[
  M(q^{m})^{m} \subset M(q)^{m} \subset M(q^{m}) \subset M(q).
\]
\note{ces deux lignes de formules sont écrites dans la moitié gauche du
feuillet, en regard du texte qui suit ; « $= p^{n}$ » est noté sous le
$q$ de $M(q)$}

Soit $M(q)$ le s-gpe de $\mathbf{Q}'(i)^{*}$ qui se définit en termes de
\ill{} ci-dessus, soit \ill{} $f_{q} \in G(q)$ \uncertain{son} dual, qui
est le gpe de Galois motivique de $\mathbf{F}_{q}$. Considérons
l'extension $\mathbf{F}_{q^{m}}$ de $\mathbf{F}_{q}$, \struck{\ill{}}

\marginal{$M(q) \simeq \mathbf{Z} \times M^{\circ}(q)$ si $q = p^{r}$,
\uncertain{$r$ pair} ; $M^{\circ}(q)$ formé des \uncertain{unités}
$u \in \mathbf{Q}'(i)^{*}$ \uncertain{exc.\ $p$}, et \ill{} \ill{}
\uncertain{conjugués entre eux}) qui satisfont à
$v_{\mathfrak{p}}(u) + v_{\mathfrak{p}'}(u) = 0$ ; le \ill{} $\mathbf{Z}$
correspond \ill{} de $q^{1/2}$. Les applications de transition entre les
$M(q)$ \uncertain{respectent} la \ill{} facteur $\mathbf{Z}$, et la
multiplication \ill{} $2$, \ill{} sur $M^{\circ}$.}
\note{note marginale écrite en oblique le long du bord gauche, d'une encre
plus bleue ; un $\mathbf{Z}$ y est surchargé, et un premier « $M^{\circ}(p)$ »
corrigé en $M^{\circ}(q)$}

\struck{$G/\ill{}$}
\[
  G(q) \longleftarrow G(q^{m}), \qquad f_{q^{m}} \longmapsto (f_{q})^{m} ;
\]
\note{sur la page, la seconde flèche est tracée comme la première, de droite
à gauche, $(f_{q})^{m}$ écrit sous $G(q)$ et $f_{q^{m}}$ sous $G(q^{m})$}
on voit que l'homom.\ transposé \uncertain{n'est autre que}
\[
  M(q) \xrightarrow{\ \text{élévation à la } m\text{-ième puissance}\ } M(q^{m})
\]
dont le noyau est \ill{} de \uncertain{(points de torsion)}
\struck{\ill{}} \uncertain{fini} (dans \ill{} $\mathbf{Z}/m\mathbf{Z}$)
et dont l'image est $M(q)^{m}$ \uncertain{contenu dans} $M(q^{m})$,
\uncertain{et contient} $M(q^{m})^{m}$, dont le \uncertain{conoyau} est
annulé par $m$. Par suite, le \struck{conoyau} noyau de \ill{} \ill{} est
annulé par $m$.

\struck{Consi} Considérons la limite \uncertain{inductive} des $M(q)$ ($q$
variant multiplicativement) par les homom.\ précédents : on trouve
$M(q^{\infty})$, un gpe \uncertain{sans torsion}, dont le dual doit être le
groupe de

\page{103}
Galois \uncertain{motivique} de $\overline{\mathbf{F}}_{p}$.
\struck{\ill{} \ill{}}
\struck{$M^{\circ}(q^{m}) \to M^{\circ}(q)$} \struck{\uncertain{induisent}}
\struck{un} \struck{\uncertain{isom.}}
\struck{$M^{\circ}(q') \otimes_{\mathbf{Z}} \mathbf{Q}$}
\struck{$\to M^{\circ}(q) \otimes_{\mathbf{Z}} \mathbf{Q}$,}
on constate que $M^{\circ}(q^{\infty})$ \struck{\ill{}} s'identifie à
\struck{\ill{}} $M^{\circ}(q) \otimes_{\mathbf{Z}} \mathbf{Q}$, qui est
lui-même \ill{} à un quotient $P/U$, $U$ étant le gpe des racines de l'$1$,
et $P$ étant le s-gpe de $\overline{\mathbf{Q}}^{*}$ formé des $\lambda$
tels que $\lambda^{n} \in M^{\circ}(q)$ pour $n$ assez grand.

\marginal{On a donc $M(q^{\infty}) \simeq \mathbf{Z} \times M^{\circ}(q^{\infty})$,
$M^{\circ}(q^{\infty}) = \varinjlim M^{\circ}(q^{m}) \simeq M^{\circ} \otimes_{\mathbf{Z}} \mathbf{Q}$}
\note{en haut à droite du feuillet}

\struck{L'homom.}
\[
  M^{\circ}(q = p^{n}) \longrightarrow P/U
\]
s'obtient par
\[
  \lambda \longmapsto \lambda^{1/n} \bmod U
\]
(à la détermination de $U$ \ill{} \ill{} près), car bien
$(\lambda^{1/n})^{n} = \lambda \in M^{\circ}(p^{n})$ ; \ill{}
$M(p^{n})^{\wedge} \subset M(p)$. \uncertain{Trouve-t-on} $P/U$ tout
entier ? $\hat{M}(p)$ \ill{}

\struck{\ill{}} \ill{} $\lambda \in \overline{\mathbf{Q}}^{*}$ tels que
$\lambda^{n} \in M^{\circ}(p)$, \struck{donc} \ill{} \struck{\ill{}}
pour \ill{} il \ill{} $\mu \in M^{\circ}(p^{m})$, i.e.\ \ill{}
$\lambda^{m} \in M^{\circ}(p^{m})$ \uncertain{pour un} ensemble ?? \ill{}
la bonne \ill{} \ill{}. On peut \struck{\ill{} $\lambda$ \ill{}} \ill{}
multiplier par $n$, et la question devient si pour $\lambda \in M^{\circ}(p)$,
$n \in \mathbf{N}^{*}$ étant donnés, il existe un $h \in \mathbf{N}^{*}$
tel que $\lambda^{h} \in M^{\circ}(p^{hn})$.
\note{à partir de « $M^{\circ}(q = p^{n}) \to P/U$ », tout le reste du
feuillet est enfermé dans un cadre (ouvert par « L'homom. » biffé) et
barré de deux longs traits obliques}

\page{105}
La forme plus juste, « géométrique » de la conjecture ci-dessus devrait
\uncertain{peut-être} \struck{aussi} bien \ill{} ainsi, \ill{} \ill{}
\ill{} : $M^{\circ} \otimes_{\mathbf{Z}} \mathbf{Q}$ (ou \ill{} $\pi$
\ill{} …), la \struck{grande} « partie compacte » du \struck{\ill{}}
groupe de Galois motivique de $\overline{\mathbf{F}}_{p}$, qui serait
\ill{} \ill{} :
\[
  G \simeq D(M^{\circ} \otimes \mathbf{Q}) \times \mathbf{G}_{m}
\]
avec l'homom.\ inj : $\mathbf{G}_{m} \to G$ s'identifiant à $j$, et
\struck{\ill{}} $(x, \lambda) \mapsto \lambda^{2}$ s'identifiant à
$\varepsilon$.

En termes de \ill{} \ill{} de Frobenius :

\marginal{\uncertain{Résulterait} \struck{\ill{}} \uncertain{probablement}
de la théorie de la multiplication complexe : cf.\ Taniyama--Shimura}

\emph{Conjecture affaiblie}. Soit $\lambda \in \overline{\mathbf{Q}}^{*}$,
\struck{tel que pour tt}
\struck{\uncertain{il existe} un entier $n > 0$ avec $\lambda^{n}$
satisfaisant aux} un \emph{entier} algébrique,
\struck{conditions, \ill{} de \ill{} $p$}
\struck{tel qu'il existe un entier $n$} \ill{} entier $i \geqslant 0$) \ill{}
tel qu'il existe un \ill{} et une \emph{unité} $u$ \ill{} pour
$\lambda = p^{i/2} u$, \struck{soit \ill{}} \ill{} et qu'il existe un
entier (\ill{} \uncertain{exc.\ $p$}]) $n > 0$ tel que
$u^{n} \in \mathbf{Q}'(i)^{*}$, et
$v_{\mathfrak{p}}(u^{n}) + v_{\mathfrak{p}'}(u^{n}) = 0$ pour tt couple
d'idéaux premiers de $\mathbf{Q}'(i)$ \uncertain{au-dessus} d'un idéal
premier de $\mathbf{Q}'$. Alors il existe un entier $m > 0$ tel que
$\lambda^{m} = (p^{m})^{i/2} u^{m}$ soit \struck{une} valeur propre du
frobenius $f_{p^{m}}$ dans un motif semi-simple \uncertain{de poids $i$}
convenable sur $\mathbf{F}_{p^{m}}$.
\note{les quatre premières lignes de l'énoncé sont en partie biffées d'un
long trait chacune, et l'état final se lit mal ; « de poids $i$ » est
ajouté dans l'interligne, au-dessus de « convenable » ; un trait vertical
isole l'énoncé à gauche}

\page{107}
$K$ extension tot.\ réelle [i.e.\ si $K$ défini par
$H \subset \pi_{1}(\mathbf{Q})$, $H$ contient tous les conjugués de $\tau$,
$\tau$ la conjugaison ordinaire)

A) $L$ extension quadratique, $L = K(x)$, $x^{2} + ax + b = 0$
($a, b \in K$), $a = x + \bar{x}$, $b = x\bar{x}$ ; $\sigma$ l'élément non
trivial de $\mathrm{Gal}(L|K)$, alors $\sigma x = \bar{x}$ pour tt $x$.
Conditions équivalentes :
\begin{quote}
$L$ tot.\ imaginaire

$\Updownarrow$

Pour tt plongement de $K$ dans $\mathbf{R}$,
$\Delta = a^{2} - 4b = (x - \bar{x})^{2}$ est $< 0$, i.e.\ $\Delta$ est
tot.\ \struck{positif} [str.]\ négatif

$\Updownarrow$ \struck{$L$ \ill{} stable par $x \mapsto \bar{x}$, \ill{}}
\struck{$(x - \bar{x})^{2}$ est tot.\ positif}
\end{quote}
\struck{Alors $\sigma$ est la conj.\ \ill{}}
\note{la signification de $a$ et $b$ (« $a = x + \bar{x}$, $b = x\bar{x}$ »,
de signe écrit tel quel) est ajoutée sous l'équation ; « str. » est
écrit au-dessus de « positif » biffé}

B) Soit $x \in \overline{\mathbf{Q}} \subset \mathbf{C}$. Conditions
équivalentes :
\begin{enumerate}
  \item[1)] $x + \bar{x}$ [et $x\bar{x}$ sont] tot.\ réels, $(x - \bar{x})^{2}$
  est tot.\ \struck{$\geqslant 0$ positif} [str.]\ $< 0$ négatif
\end{enumerate}
\marginal{[\uncertain{conséquence} des autres, car
$x\bar{x} = \frac{1}{4}\bigl((x + \bar{x})^{2} - (x - \bar{x})^{2}\bigr)$]}
\note{la note marginale est reliée par une flèche à « [et $x\bar{x}$
sont] »}

$\Updownarrow$
\begin{enumerate}
  \item[2)] \struck{$x$ est quadr.}\ $\exists$ un $K$ tot.\ réel et une
  extension $L$ tot.\ imaginaire de $K$, dont $x$ soit un générateur.
\end{enumerate}

$\Updownarrow$
\begin{enumerate}
  \item[2')] $x \in \mathbf{Q}'(i) - \mathbf{Q}'$, i.e.\ \ill{}
\end{enumerate}

En effet, 2) $\Rightarrow$ 1) en vertu de A), et pour 1) $\Rightarrow$ 2),
il suffit de prendre $K = \mathbf{Q}(a, b)$ où $a = x + \bar{x}$,
$b = x\bar{x}$, et appliquer encore A). \emph{Remarque}. Cet
\uncertain{argument} prouve que les conditions de 1) sont stables par
conjugaison, ce qui n'était pas évident autrement.

\struck{3) $x\bar{x}$ est tot.\ str.\ positif, et de même pour
$x_{\alpha}\bar{x}_{\alpha}$ pour tt conjugué de $x$}

C) Soit $x \in \overline{\mathbf{Q}} \subset \mathbf{C}$. Conditions
équivalentes :
\begin{enumerate}
  \item[1)] $x + \bar{x}$ tot.\ réel, $(x - \bar{x})^{2}$ tot.\ négatif
  \item[2)] $\exists$ $K$, $L$ comme plus haut, $x \in$ \struck{$K$}
  \uncertain{$L$}
  \item[3)] $x \in \mathbf{Q}'(i)$.
  \item[4)] $x\bar{x}$ tot.\ positif, \struck{et de même} ainsi que
  $x_{\alpha}\bar{x}_{\alpha}$ pour tt conjugué $x_{\alpha}$ de $x$.
\end{enumerate}
\note{des doubles flèches relient 1), 2), 3) ; une ligne partie de 3)
va vers 4), marquée d'une flèche « ? $\Uparrow$ » ; un « égal » est biffé
au-dessus de « ainsi que »}

\page{109}
\note{à gauche, en tête, un diagramme d'extensions, redessiné ici :
$\mathbf{Q}'(i)$ — $L$ en haut, $\mathbf{Q}'$ — $K$ au-dessous, $L$ au-dessus
de $K$ et de $L_{0}$, $L_{0}$ au-dessus de $K_{0}$, $K$ et $K_{0}$ au-dessus
de $\mathbf{Q}$ ; une accolade, à droite de $L$, $L_{0}$, $K$, $K_{0}$, porte
« galoisienne »}

\begin{tikzcd}[column sep=small, row sep=small]
  \mathbf{Q}'(i) \arrow[r, no head] \arrow[d, no head] & L \arrow[d, no head] \arrow[dr, no head] & \\
  \mathbf{Q}' \arrow[r, no head] & K \arrow[dd, no head] \arrow[dr, no head] & L_{0} \arrow[d, no head] \\
  & & K_{0} \\
  & \mathbf{Q} &
\end{tikzcd}

On veut construire le gpe de type \uncertain{multiplicatif} $G(q)$
\uncertain{associé} à l'ens.\ $M(q)$ des $\lambda \in L$ tels que
$\lambda\bar{\lambda} = q^{i}$, et $\lambda$ entiers exc.\ $p$. On a vu que
\[
  M_{0}(q)/\mathrm{Torsion} \simeq (\mathbf{Z}^{\mathcal{P}})' \qquad (*)
\]
\marginal{N.B. $M_{0}(q) = M_{0}(p)$ = ens.\ des unités $u \in L$ entières
exc.\ $p$ et satisfaisant $u\bar{u} = 1$.}

$\mathcal{P}$ l'ens.\ des idéaux premiers de $L$ sur $p$. On peut supposer
que les idéaux \struck{pre} de $K$ se \uncertain{décomposent} dans $L$.
Choisissons un des \struck{\ill{}} idéaux \ill{} $\mathfrak{p}$, on peut
\uncertain{définir} \ill{} sur $M$ une forme \ill{}
$\lambda \mapsto v_{\mathfrak{p}}(\lambda)$ \ill{} $\check{M}$ \ill{}
\ill{} $\mathfrak{G}$ le \ill{} conjugué d'ailleurs \ill{} gpe de Galois
$\Gamma$ de $L$, d'où
\[
  \mathbf{G}_{m\,L} \xrightarrow{\ v_{\mathfrak{p}}\ } G_{L}.
\]
Soit $\Gamma_{\mathfrak{p}}$ le gpe de décomposition pour $\mathfrak{p}$,
\uncertain{quotient} \struck{$K_{0}$ le corps} $L_{0}$ le corps
\struck{aussi} des invariants de $\Gamma_{\mathfrak{p}}$, $K_{0} = L_{0} \cap K$.
Alors $L_{0}$ est le plus petit corps de rationalité de $v$, \uncertain{sur}
lequel \struck{dans}
\[
  v : \mathbf{G}_{m\,L_{0}} \longrightarrow G_{L_{0}}
\]
est induit par $v_{\mathfrak{p}}$ sur $L_{0}$.

D'autre part, si $\mathfrak{p}_{0}$ est induit par $\mathfrak{p}$,
l'\ill{} local complété \struck{l'anneau \ill{} \ill{}} alors
\struck{\ill{} \ill{} \ill{} \ill{}} \ill{} de $\mathfrak{p}_{0}$ et
\ill{} le corps $(L_{0}, \mathfrak{p}_{0})$ est défini \ill{}
$\mathbf{Z}_{\ell}$. Notons que \ill{} la donnée de la classe à isom.\
unique près par la donnée de la \struck{\ill{}} \ill{} de $L$, et par sa
définition \ill{} relatif de $\mathbf{Z}_{(\mathfrak{p})}$ dans $L$ ….
Notons que $L_{0}$ est \ill{} imaginaire, grâce à l'hypothèse que
$\tau\mathfrak{p} \neq \mathfrak{p}$.

\page{111}
\[
  \struck{\ill{}}\ \mathbf{Z}/2\mathbf{Z} \to \Gamma \to \Gamma' \to 0,
  \qquad 1 \to \Gamma_{\mathfrak{p}} \subset \Gamma
\]
\note{en haut à gauche ; $\Gamma_{\mathfrak{p}}$ est écrit sous $\Gamma$,
avec un signe d'inclusion vertical, et un trait vertical relie
$\mathbf{Z}/2\mathbf{Z}$ au « 1 »}

D'autre part, $(\mathbf{Z}^{\mathcal{P}})'$ correspond à un tore sur
$\mathbf{Q}$, qui peut se décrire ainsi, comme \ill{} \ill{} \ill{} en
termes de Galois : l'ext.\ quadratique $L_{0}$ de $K_{0}$ définit un
\uncertain{groupe tordu}
$U$ de \ill{} (savoir
$\operatorname{Ker}\bigl(\prod_{L_{0}/K_{0}} \mathbf{G}_{m\,L_{0}} \to \mathbf{G}_{m\,K_{0}}\bigr)$)
et on considère \ill{}
\[
  G' = \prod_{K_{0}/\mathbf{Q}} U = \operatorname{Ker}\Bigl(\prod_{L/\mathbf{Q}} \mathbf{G}_{m\,L_{0}} \to \prod_{K_{0}/\mathbf{Q}} \mathbf{G}_{m\,K_{0}}\Bigr)
\]
(donc
$G'(\mathbf{Q}) = \operatorname{Ker}\bigl(L_{0}^{*} \xrightarrow{N_{L_{0}/K_{0}}} K_{0}^{*}\bigr)$).
\ill{} \ill{} un homom.
\note{sous le premier produit du noyau, l'indice est lu « $L/\mathbf{Q}$ »,
non « $L_{0}/\mathbf{Q}$ »}
\[
  (**) \qquad \varphi : G' \longrightarrow G
\]
qui est une isogénie de $G'$ sur $G^{\circ}$. D'ailleurs
$v_{\mathfrak{p}} : \mathbf{G}_{m\,L_{0}} \to G_{L_{0}}$ se factorise par
$G'_{L_{0}}$ ; d'où

\begin{tikzcd}
  \mathbf{G}_{m\,L_{0}} \arrow[r, "v'"] \arrow[d, no head, "\Vert" description] & G'_{L_{0}} \arrow[d] \\
  \mathbf{G}_{m\,L_{0}} \arrow[r, "v"] & G_{L_{0}}
\end{tikzcd}

On fera attention que l'élément $f_{0}$ de \uncertain{$G^{\circ}(\mathbf{Q})$}
\marginal{$= f\, j(p^{1/2})$}
ne se relève en un élément $f'$ de $G'(\mathbf{Q})$ que si
\struck{\ill{}} \ill{} \ill{} i.e.\
$M(q)/\mathrm{torsion} \xrightarrow{\ \sim\ } (\mathbf{Z}^{\mathcal{P}})'$.
\note{« $= f\,j(p^{1/2})$ » est écrit au-dessus de $G^{\circ}(\mathbf{Q})$,
dans un contour relié par un trait à $f_{0}$ ; l'exposant de $G$ est mal
formé, et un astérisque suit la parenthèse}

\page{113}
\marginal{demander à Serre ?}
Ce n'est \ill{} \ill{} que \ill{} le cas, \struck{\ill{}} si
$K_{0} = \mathbf{Q}$ : il s'agit de donner un exemple de
$\mathbf{Z}/2\mathbf{Z}$ opérant non trivialement sur $\mathrm{Pic}(A)$, $A$
étant \struck{l'anneau} \struck{\ill{}} des entiers de $L$. On constate
d'ailleurs que $\mathbf{Z}/2\mathbf{Z}$ opère sur $\mathrm{Pic}(A)$ par
\uncertain{symétrie} ; il s'agit donc de trouver $\mathrm{Pic}(A)$ non un
$2$-gpe !

\page{116}
Soit $G$ un gpe \struck{\ill{}} alg.\ \ill{} sur $\mathbf{Q}$ provenant
d'une situation motivique \uncertain{[$\mathcal{M}$]}, \struck{\ill{}}
soit $U = \operatorname{Ker}(G^{\circ} \xrightarrow{\ \varepsilon\ } \mathbf{G}_{m})$.
On a donc un élément canonique \ill{} $\mathbf{G}_{m}$.
\struck{On a $U = \ill{}\ G$}

\marginal{N.B. La \uncertain{condition} \ill{} \ill{} \ill{}
\uncertain{structure} \ill{} 1, 2, 3}

\emph{Proposition 1}. Conditions équivalentes :
\begin{enumerate}
  \item[(i)] $U$ est un tore.
  \item[(ii)] $\varepsilon$ est \ill{} dans l'ens.\ des \uncertain{caractères}
  $G \to \mathbf{G}_{m}$ (« normalisés »)
  \item[(iii)] $\mathcal{M}_{\mathbf{R}}$ admet un foncteur fibre (« normalisé »),
  i.e.\ $\mathcal{G}_{\mathbf{R}} \neq \emptyset$.
\end{enumerate}
C'est \uncertain{toujours} vrai si $\mathcal{M}$ a degrés pairs, i.e.\
\uncertain{$j(\mu_{2}) = 0$}.
\note{les deux « normalisé » sont ajoutés au crayon ; le premier, entre
parenthèses, est relié par une accolade au foncteur fibre de (iii)}

\marginal{i.e.\ correspondant à un \ill{} \ill{} \ill{} \ill{} \ill{}
de foncteurs \ill{} \ill{} dont \ill{} \ill{}}
\note{note marginale au crayon, dans un contour, le long du bord gauche}

\emph{Proposition 2}. Conditions équivalentes :
\begin{enumerate}
  \item[(i)] $U$ est de la forme $U^{\circ} \times \mu_{2}$, $U^{\circ}$ un tore
  [dans $\mathrm{Hom}(\mathcal{L}; \mathbf{G}_{m})$]
  \item[(ii)] $\varepsilon$ n'est \struck{\ill{}} pas un \uncertain{carré}
  « normalisé »
  \item[(iii)] $\mathcal{M}_{\mathbf{R}}$ n'admet pas de foncteur fibre, i.e.\
  $\mathcal{G}_{\mathbf{R}} = \emptyset$.
\end{enumerate}

Le fait que (i) implique \ill{} ; (ii) [dans prop.\ 1 et prop.\ 2,
résulte trivialement] de \struck{l'équivalence de} l'existence de $j$ avec
$\varepsilon j(\lambda) = \lambda^{2}$ ; ainsi que le fait que
$(1\,\mathrm{ii}) \Longleftrightarrow (\text{non}\ \struck{\ill{}}\ 2\,(\mathrm{ii}))$.

On a (iii) $\Rightarrow$ (ii), car plus généralement pour tt
\struck{\ill{}} objet de [degré impair] \struck{degré} de
$\mathcal{M}_{\mathbf{R}}$, le \uncertain{rang} \struck{degré} doit être
pair, comme il résulte de l'existence d'une forme alternée fondamentale.
D'autre part (i) $\Rightarrow$ (iii), car si $T$ est un tore
\uncertain{compact} sur $\mathbf{R}$, \ill{}
$H^{2}(\mathbf{R}, T) = 0$ ! Enfin, si \struck{Cor} $\mathcal{M}$ est à
degrés pairs i.e.\ $j(\mu_{2}) = 0$, alors $j = 2j'$ et
$\varepsilon j' = \mathrm{id}$, donc on a (ii).
\note{l'exposant de $H^{2}(\mathbf{R}, T)$ est douteux}

\emph{Cor}. Si $\mathcal{M}$ est quelconque, la sous-catégorie [motivique]
$\mathcal{M}^{\mathrm{pair}}$ \struck{\ill{}} (qui correspond à
$G/\struck{\ill{}}\,j(\mu_{2})$) \struck{\ill{}} admet un
$\mathbf{R}$-foncteur fibre.

\page{118}
\emph{Cor}. Dans le cas de la prop.\ 2, \struck{la catégorie} \ill{}
$G^{\circ} \simeq U^{\circ} \times \mathbf{G}_{m}$, $\varepsilon = \mathrm{pr}_{2}$,
$j(\lambda) = \struck{\ill{}}$ \uncertain{$j_{1}(\lambda)^{2}$}. La structure
\struck{motivique 1) 2) 3)} de catégorie tensorielle de
$\mathcal{M}_{\mathbf{R}}$ est connue tot.\ \ill{} par la connaissance de
$G$ et l'élément de $H^{2}(\mathbf{R}, U)$ caractéristique, qui
\struck{\ill{}} \ill{} l'unique élément non nul de $H^{2}(\mathbf{R}, U)$,
correspondant à l'unique élément non nul de $H^{2}(\mathbf{R}, \mu_{2})$.
Il existe un \struck{\ill{}} objet $E$ de $\mathcal{M}_{\mathbf{R}}$
\struck{\ill{}} de degré $1$, rang $2$, \struck{\ill{}} tel que la
sous-catégorie tensorielle \struck{\ill{}} $\mathcal{M}'$ qu'il engendre
\ill{} \ill{} $\mathbf{R}$-tordue, et \struck{\ill{} \ill{}} tel que
$U^{\circ}$ y opère trivialement.

Dans le cas favorable $\mathcal{G}_{\mathbf{R}} \neq \emptyset$,
\struck{i.e.\ $U$ un tore \ill{}} l'ens.\ [$H^{0}(\mathcal{G}_{\mathbf{R}})$
des objets à isom.\ près de $\mathcal{G}_{\mathbf{R}}$]
\struck{$\mathrm{Ob}\,\mathcal{G}_{\mathbf{R}}$} est un \struck{ens.\
principal homogène} torseur sous le gpe $H^{1}(\mathbf{R}, U)$, qui
lui-même \ill{} \ill{}
${}_{2}U(\mathbf{R}) \simeq (\mathbf{Z}/2\mathbf{Z})^{\dim U}$.
D'autre part, à chaque élément de $\mathrm{Ob}\,\mathcal{G}_{\mathbf{R}}$
correspond une \struck{structure de \ill{}} [structure polarisante] sur
$G^{\circ}$, définie par un homom.\
$\mu_{2\,\mathbf{C}} \xrightarrow{\ i_{1}\ } G^{\circ}_{\mathbf{C}}$ tel
que $\varepsilon_{\mathbf{R}} i_{1} = \mathrm{id}$,
$i_{1}\bar{\imath}_{1} = j_{\mathbf{R}}$, et l'ens.\ de ces structures est
\ill{} un torseur sous ${}_{2}U(\mathbf{R})$. L'application
\[
  H^{0}(\mathcal{G}_{\mathbf{R}}) \longrightarrow \mathrm{Pol}(\mathcal{G}_{\mathbf{R}})
\]
est compatible avec les opérations de ${}_{2}U(\mathbf{R})$, \emph{donc
bijective}. Donc les $\mathbf{R}$-foncteurs fibres à isom.\ près

\page{120}
et les structures de polarisation sur $G^{\circ}$ se déterminent
\emph{mutuellement}.

\struck{Donc} [À voir : \uncertain{le gpe des automorphismes} de
$(G^{\circ}_{\mathbf{R}}, \varepsilon_{\mathbf{R}}, j_{\mathbf{R}})$
voire $(G^{\circ}, \varepsilon, \mu)$, voire $(G, \varepsilon, \mu)$
opère-t-il transitivement sur $\mathrm{Pol}(G)$ ? \struck{Et de même}
\uncertain{Par exemple}, est-il impossible \struck{\ill{}} de distinguer
\struck{\ill{}} une \struck{\ill{}, \ill{}} polarisation ou une
\uncertain{classe} de polarisations préférentielles ?]
\note{l'indice de $G^{\circ}$ dans le deuxième triplet est mal formé}

Prenons p.\ ex.\ une catégorie de Motifs géométriques sur $\mathbf{F}_{q}$,
$q = p^{n}$. Alors \ill{} d'ailleurs il y a une classe de conjugaison
d'homom.
\[
  i : \mathbf{G}_{m\,\mathbf{C}} \longrightarrow G_{\mathbf{C}}
\]
distinguée, avec
\[
  \bar{\imath}\, i = j, \qquad \varepsilon i = \mathrm{id}
\]
définie, \struck{\ill{}} si $G$ correspond à $M \subset L^{*}$ ($L$
quadratique tot.\ imag.\ sur $K$ tot.\ réel), par la \uncertain{forme}
$\lambda \mapsto \frac{1}{n} v_{\mathfrak{p}}(\lambda)$, où $\mathfrak{p}$
est un idéal premier de $L$ sur $p$.
[$v_{\mathfrak{p}}(\lambda)$ \add{$+$} $v_{\mathfrak{p}}(\bar{\lambda}) =
v_{\mathfrak{p}}(\lambda\bar{\lambda}) = v_{\mathfrak{p}}(q^{j^{*}(\lambda)}) = j^{*}(\lambda)$ ;
$v_{\mathfrak{p}}(q) = 1$]. En fait, il convient de « normaliser »
$G_{\mathbf{C}}$ [par le] choix (sur $\mathbf{Q}$ !) d'un tel $i$, (NB. Un
autom.\ de $G^{\circ}$ qui \struck{induit} respecte $i$ est l'identité !).
La structure d'\uncertain{effectivité} sur $\mathcal{M}$ est déterminée
grâce à $i$, \struck{elle} sur $\mathcal{M}_{K}$ ($K$ ext.\ de $\mathbf{Q}$)
en particulier de $\mathcal{M}_{\mathbf{R}}$, $\mathcal{M}_{\mathbf{C}}$,
par la connaissance du système des \ill{} conjugués
\note{la phrase se poursuit au feuillet 122 (lot 7), qui s'ouvre sur « de
$i$ ». La forme « $\frac{1}{n} v_{\mathfrak{p}}$ » porte « $\frac{1}{n}$ »
entouré et relié par un trait à une note marginale cerclée}

\marginal{N.B. $i$ est \ill{} \uncertain{défini} sur $\mathbf{Q}$, si
$[K : \mathbf{Q}] = \ill{}$ (car \struck{\ill{}} \ill{}
$\mathrm{Hom}(\mathbf{G}_{m\,\mathbf{C}}, G^{\circ}_{\mathbf{C}})$)}
\marginal{\struck{cet} \uncertain{Calculer} cette forme \struck{\ill{}}
\ill{} par degrés : valeurs entières \ill{} \ill{} s-groupe \ill{}}
\marginal{\ill{} $\mathcal{M}$ où $j = 0$ ; ex : prendre $q = p \neq 2$,
$L = \mathbf{Q}(\sqrt{-p})$, tel que nb de classes de $L$ soit $1$, p.\ ex.\
$p = 3, 7, 11$, …}
\note{trois notes marginales en oblique le long du bord gauche, les deux
dernières dans des contours au crayon ; le « $\frac{1}{n}$ » cerclé du
texte renvoie à la deuxième}

\end{document}
